\section{The budget rule, pre-registration hashes and amendments}
\label{sec:hashes}

\subsection{The budget rule}
\label{sec:budget-detail}

We fixed the budget of every arm except Baseline A on a development run on fold 0: BNM's by a
rule declared before that run, which reads its nuclear norm (the $\gamma$ sweep on the same fold
had already indicated that the rule would give about \bnmBudgetPreview{} steps, and it gave
\bnmBudgetSteps{}); the others with its TPR curve in view. Baseline A's 8000 steps were fixed before any results existed. We read each arm at its
last checkpoint, so the training budget is the model selector. Every confirmatory run
reached its arm's final budget. BNM's ran to 4000 steps so that all its readouts came from one
run; no part of its training depends on the total number of steps, so its step-1000 checkpoint
is the model a 1000-step run would produce. For some arms we stated the budget as a rule derived
from the arm's configuration: 4000 for DANN, twice its $\lambda$
horizon; 4000 for the $\lambda=0$ control, read where its treatment is read; and 1000 for BNM,
twice the step at which its nuclear norm reaches an analytic ceiling. Baseline A uses 8000, the
matched budget of the frozen comparison. For CORAL we declared 2000 on cost, declining the
development curve's maximum at 2500. For FreeMatch we first declared 2000 on robustness
grounds, with the development curve in view, and later amended it to 1500, the development
curve's highest point, although Section~\ref{sec:negative} shows that this maximum is poorly
located. Only \rwBudgetPreCurve{} were fixed before the corresponding
development curve existed; the others were chosen with that curve in view. BNM's 1000 matches
no other arm's budget, so we read it at 4000, the step of its control. Where we compare two arms
at the same step, a budget match removes only the budget confound; the head and tuning
differences of Section~\ref{sec:arms} remain.

\subsection{Pre-registration}
\label{sec:prereg-detail}

We fixed each arm's budget and read-off step before its own confirmatory folds ran and, except
for the four items below, hash-protected them then. The $\lambda=0$ control's folds 1--4 and the
four paired contrasts against it came later (Section~\ref{sec:prereg}): on 2026-09-01 we recorded
in the DANN pre-registration that the control would be reported from the development fold alone
(key \texttt{arm1\_folds\_1\_4\_declined}), and on 2026-09-06, after the FreeMatch, CORAL and
DANN confirmatory folds had run, we reversed that decision. The control's own pre-registration
precedes all its confirmatory folds, but we chose its contrasts knowing the adaptation arms'
results.
Two versions of our private Kaggle dataset hold copies of these files. Their dates are the
version dates shown on the dataset page, which we read ourselves; a reader cannot check them, so
they corroborate our account but are not independent timestamps. One holds the DANN declaration
(see below). The other, dated \lambdaKaggleDate{}, before the $\lambda = 0$ control's first
confirmatory fold, holds that control's pre-registration, which fixes the four paired contrasts
against it, and the pre-registration's audit; both match their recorded digests. Four items rest on our run logs alone, which are not distributed (Data and code
availability): Baseline A's budget, fixed
before any results in our first frozen comparison, which predates our hash protection; the
precedence of one DANN prediction over its development-fold test; the timing of BNM's budget;
and FreeMatch's amended budget, whose hash file did not survive. We evaluated every condition
declared for the $\lambda=0$ control at step zero, and across its whole development run, before
writing it into the pre-registration. Fold 0 enters no confirmatory number: every contrast and
every unfrozen-arm mean uses folds 1--4, 391 test positives in all.

Table~\ref{tab:hashes6} lists every pre-registration and budget file with its digest, re-derived
from disk when the paper is built. We record amendments beside a hashed file rather than
rewriting it.

% GENERATED FILE -- do not edit by hand.
% Rebuild:  python make_all.py (generator: make_dilute10_tables.py)
% Appendix B (paper v6): tables/A2_hashes_v2.tex unchanged, plus the 1:10 experiment's declarations, digests re-derived from disk through scripts/integrity.py.
\begin{table}[htbp]
\centering
\caption{\textbf{Protected artifacts and when each was recorded.} Each enrolled file is made read-only once written and is checked by digest at every build. \emph{Recorded} is the file's own timestamp; where it has none, \emph{Source} says whether its timing rests on a flag inside the file, on a dated Kaggle dataset version, or on our run logs alone. \emph{vs.\ runs} compares it with the start of the first confirmatory run it governs. Two kinds of digest are checked: a whole-file digest, and a self-certifying digest of the file body with its own hash field removed. \emph{Enrolled} says whether the file is in the inventory that every analysis run re-checks; every row's digest, enrolled or not, matches its file when the table is built. Fourteen artifacts are hashed; thirteen are re-checked by every analysis run, and the window-averaging declaration is hashed but not enrolled. The 1:10 rows are dated against that experiment's first run. The DANN pre-registration has two rows, its archived pre-fold copy and the file as regenerated after the folds. A date without a time is a day-only record. The last row is protected but governs no run.}
\label{tab:hashes6}
\scriptsize
\setlength{\tabcolsep}{4pt}
\begin{tabular}{>{\TRAG}p{2.6cm}>{\TRAG}p{1.3cm}>{\TRAG}p{1.75cm}>{\TRAG}p{1.3cm}>{\TRAG}p{1.3cm}>{\TRAG}p{1.4cm}l>{\TRAG}p{1.4cm}}
\toprule
Artifact & Recorded (UTC) & Source & First run (UTC) & vs.\ runs & Digest & sha256 (first 16) & Enrolled \\
\midrule
BNM budget & --- & run logs only & 2026-09-14 11:36 & not checkable & whole file & \texttt{5b92e161ed1e301f} & yes \\
BNM pre-registration & --- & file (flag only) & 2026-09-14 11:36 & before (flag) & whole file & \texttt{d205f6784213b75d} & yes \\
CORAL budget & 2026-08-21 18:01 & file & 2026-08-23 12:42 & before & self-certifying & \texttt{65c20810485968d8} & yes \\
DANN budget + sensitivity readout & 2026-09-02 & file & 2026-09-03 13:54 & before & self-certifying & \texttt{abf0dae0e313ae53} & yes \\
DANN pre-registration, pre-fold copy & 2026-09-02 16:26 & archived copy & 2026-09-03 13:54 & before & self-certifying & \texttt{51351308196ad332} & no (archived) \\
DANN pre-registration, regenerated file & 2026-09-05 11:43 & file & 2026-09-03 13:54 & \textbf{after} & self-certifying & \texttt{3ad1fc361a150b95} & yes \\
$\lambda=0$ control pre-registration & --- & Kaggle dataset version, 6 Sept 2026 & 2026-09-07 08:38 & before & self-certifying & \texttt{c1694a479eafb040} & yes \\
Checkpoint-window averaging declaration & --- & file (flag only) & --- & before (flag) & self-certifying & \texttt{a572b81e4c730590} & \textbf{no} \\
\midrule
\multicolumn{8}{l}{\emph{The 1:10 diluted-target experiment (Section~\ref{sec:dilute-methods}):}} \\
1:10 pool definition & --- & digest in part 1 & 2026-10-13 18:58 & before & whole file & \texttt{18c73d6c1cbad99e} & yes \\
1:10 pre-registration, part 1 & 2026-10-13 13:19 & file & 2026-10-13 18:58 & before & self-certifying & \texttt{bd537a025e81833e} & yes \\
1:10 pre-registration, part 2 & 2026-10-13 18:07 & file & 2026-10-13 18:58 & before & self-certifying & \texttt{b8e09197a204be83} & yes \\
1:10 pre-registration, part 3 & 2026-10-19 14:13 & file & 2026-10-20 16:10 & before & self-certifying & \texttt{1b583ea2ae5056ed} & yes \\
\multicolumn{8}{l}{\emph{Analysis declarations (govern analyses, not runs):}} \\
1:10 analysis plan & 2026-10-23 18:45 & file & --- & n/a & self-certifying & \texttt{68367e25c784334c} & yes \\
DANN $-$ $\lambda=0$ memorisation plan (retrospective) & --- & run logs only & --- & n/a & self-certifying & \texttt{9dee7d6a02c17be0} & yes \\
\midrule
\multicolumn{8}{l}{\emph{Protected, not a pre-registration (governs no run):}} \\
BNM step-0 gradient norms, re-measured & --- & run logs only & --- & n/a & whole file & \texttt{04930f2fe924b99a} & yes \\
\bottomrule
\end{tabular}
\end{table}

\begingroup\footnotesize\noindent\textbf{Notes to Table~\ref{tab:hashes6}.} What each file fixed: 
\emph{BNM budget}: budget 1000, BNM's rule-derived budget. The file carries no timestamp and no before-flag, so its timing rests on our run logs alone.
\emph{BNM pre-registration}: the weight $\gamma$. BNM's weight came from a development-fold sweep over a grid set by gradient norms at step zero; the sweep read the nuclear norm and the predicted-positive fraction, never TPR. The budget was left undecided, because its rule reads a crossing only the development run produces.
\emph{CORAL budget}: budget 2000, declared on cost, declining the development curve's maximum.
\emph{DANN budget and sensitivity readout}: budget 4000, $2 \times$ the $\lambda$ horizon, and the step-2000 sensitivity declaration, which we call the headline-and-sensitivity rule: the headline is the mean at 4000, and the mean at step 2000 is a budget-sensitivity check reported beside the headline, not instead of it. Every confirmatory fold records this file's self-certifying digest when it starts, so the rule is verified as fixed before the folds. The declaration says ``5-fold mean'', but the same file excludes fold 0 from every headline, and we report folds 1--4; with the development fold, the five-fold means on the cleaned pool are \dDannFiveFoldFourK{} at 4000 and \dDannFiveFoldTwoK{} at 2000.
\emph{DANN pre-registration}: the arms, the abort conditions, the rule for declaring the margin $M$, the budget decision, the headline-and-sensitivity rule, and the predictions, among them the arc-suppression prediction and the spread prediction. We regenerated the file after the folds; the first row is the copy that predates them, and every one of its \dannPrefoldIdenticalKeys{} pre-fold keys is unchanged in the second. The copy is entry \dannPrefoldEntry{} of the file's hash history, and we archive it in the repository; the regeneration, on 2026-09-05, added \dannPrefoldAddedKeys{} keys reporting fold results. The copy does not establish that the arc-suppression prediction preceded its own development-fold test, whose outcome it already contains; entry \dannArcEntry{} of the hash history was not recovered.
\emph{$\lambda=0$ control pre-registration}: budget 4000, under the rule that the control is read where its treatment is read; three outcome branches (low, near the adaptation arms, in between), which decided whether BNM would run, with no run by default; and the four paired contrasts against the control. The file's description of two other arms' budgets is wrong: it attributes FreeMatch's 1500 to the saturation of its adaptive threshold, which completes 90\% of its rise on the development fold only at step \fmTauNinetyStep{}, and CORAL's 2000 to its configuration. Appendix~\ref{sec:budget-detail} gives the bases recorded at the time (FreeMatch: the development curve's highest point; CORAL: cost). The sentence governs nothing in the control's run.
\emph{Checkpoint-window averaging declaration}: $k=8$ primary, the grid $\{1,2,4,8\}$, both score combinations, and a window that ends at each arm's read-off step; it governs analyses, not training runs.
\emph{BNM step-0 gradient norms, re-measured}: a second, later measurement. The $\gamma$ grid is never re-derived from it, and its original measuring run is lost. Protected, but not a pre-registration: it governs no run.
A further digest in our run logs, \texttt{7356acda\ldots}, is the hash of a superseded pre-registration text; the text is lost, and we keep its digest in \texttt{public\_inputs/\allowbreak record\_extract.json}. The window-averaging row, marked \emph{no} under \emph{Enrolled}, is hashed in our run logs, and its digest matches the file, but it is absent from the inventory that every analysis run re-checks, so it has neither the re-check nor the read-only protection of the other thirteen. We did not enrol it after the fact, because that would date its protection to the writing of this paper. The first 1:10 run was CORAL's frozen stage on fold 1, run locally.\par\endgroup


\paragraph{The DANN pre-registration.} We recorded it in our run logs on 2026-08-26 and amended
it there on 2026-08-28 and 2026-08-30. Its machine-readable file is regenerated whole by its
writer and was last written on 2026-09-05, after DANN's confirmatory folds, so that file's own
hash dates only the regeneration. A copy of the file as it stood on 2026-09-02, before the
folds, survives in a private Kaggle dataset version dated 2 September 2026, and we archive it in
the repository. Its digest, \texttt{\dannPrefoldSha{}\ldots}, is entry \dannPrefoldEntry{} of the file's hash
history, its
recorded time precedes the first confirmatory fold, and every one of its
\dannPrefoldIdenticalKeys{} pre-fold keys is identical in the current file; the
regeneration added \dannPrefoldAddedKeys{} keys reporting fold results. This copy therefore
verifies, before the confirmatory folds, the spread prediction, the headline and sensitivity
rule, the budget, the abort conditions as we re-specified them on 2026-08-30 from the fold-0
control and development runs (the file labels that change derived, not pre-registered; it made
two conditions advisory, so it could stop no run), and every other key. It does not verify two things. The first is the precedence of the
arc-suppression prediction over its own development-fold test: the prediction is dated
2026-08-30 in our run logs, but the 2 September copy already contains its outcome, resolved on
2026-09-01, and the earlier file state (history entry \dannArcEntry{}) was not recovered; the
prediction was falsified. The second is the content of the file's states before 2 September,
history entries 1--9. The step-2000 sensitivity rule is also verifiable through the separate
DANN budget file, which predates the folds and whose digest every confirmatory fold recorded
when it started.

\paragraph{The FreeMatch budget amendment.} We amended FreeMatch's budget from 2000 to 1500
steps, the development curve's highest point, although Section~\ref{sec:negative} shows that
this maximum is poorly located; every FreeMatch confirmatory result we report was run
at 1500. Our run logs hold the original text's digest,
\texttt{7356acda\ldots}; the file that held the amended text's digest is not in the repository
and cannot be recreated, so no surviving hash covers the amended text.

\subsection{The pre-registration of the dilution experiment}
\label{sec:prereg-dilute}

The 1:10 experiment of Section~\ref{sec:dilute-methods} was declared in three parts, each
self-certifying, plus a pool file and an analysis plan; Table~\ref{tab:hashes6} lists them with
their first governed run, read from the run records.

\paragraph{Part 1} (\texttt{\dlShaPartOne\ldots}) fixed the pool and the sampler; the arms,
folds, budgets and read-off steps; the reuse of the $\lambda=0$ control and the rerun of the
CORAL and FreeMatch frozen stages; the contrasts $G$ and $D$; the three branches with their
sentences and the qualifier; the rule that the abstract and the conclusion report every arm's
branch, leading with DANN and BNM; and our \dlNPredWord{} predictions. The pool file's digest is
recorded inside it.

\paragraph{Part 2} (\texttt{\dlShaPartTwo\ldots}) was recorded after step-zero measurements of
at most \dlPartTwoMaxStepsWord{} training steps per fold, but before any confirmatory run. It
fixed the step-zero thresholds, the
tolerance band on the lens fraction of the target batches, and which inherited stop conditions
abort a run and which only log; CORAL's feature-norm guard is among the latter. DANN's four
diluted folds ran from the first version of our run scripts, which did not record part 2's
digest in the run files; our pre-run check verified part 2 before each of those runs, and every
later run records it. Part 1 states that part 2 would record the step-zero gate values of the
new CORAL and FreeMatch frozen heads. Those heads did not exist when part 2 was recorded, so
part 2 fixed the rule instead: each unfrozen run is gated against the values its frozen stage
wrote at the end of its run, in that stage's \texttt{gate\_values.json}, which we hashed before
the unfrozen run started.

\paragraph{Part 3} (\texttt{\dlShaPartThree\ldots}) was written after the first attempt of
CORAL's unfrozen run on fold 1, on \dlAbortDate{}, stopped at step 1 on that guard and produced
no data, and before every confirmatory CORAL and FreeMatch unfrozen run. It binds those two arms
only. It permits that fold's rerun, keeps CORAL as pre-registered with the guard advisory, adds
the rule that withholds the verdict from an arm whose run applies fewer than half its budget's
optimiser updates, which we set after a label-free \dlCoralCheckSteps{}-step check in which CORAL
applied \dlCoralCheckAppliedMin{}--\dlCoralCheckAppliedMax{} of \dlCoralCheckSteps{} updates on
every fold, and which fired on no run; and it requires a pre-run check of FreeMatch's early
updates before its runs. It changes nothing about how any arm is trained. Part 2 described
itself as final; part 3 adds to it without changing any earlier key.

\paragraph{The analysis plan} (\texttt{\dlShaPlan\ldots}) was hashed before any value of the
experiment was computed. Its driver was edited once afterwards, to fix an index error in the
scoring of the predictions that swapped BNM's and CORAL's $D$; every contrast and branch is
unaffected.

\paragraph{The memorisation plan} (\texttt{\dlShaGapPlan\ldots}) is retrospective and in no
pre-registration: we decided on the DANN $-$ $\lambda=0$ analysis of
Section~\ref{sec:negative} after every result was known, and declared and hashed its rule before
computing any value.
